\section{Introduction}
A payment can arrive in a Colombian account before the person for whom the money was intended can use it. In accounts of sending money to Venezuela, women describe making a local transfer, relying on a contact for conversion or delivery, and checking separately whether the recipient got the funds. This sequence illustrates an arrangement described in the migrant synthesis; it is not a reconstruction of one woman's transaction (Section~\ref{sec:migrant-journeys}). The local payment counts as digital use even though it says little about who can instruct the next step or seek help if the money does not arrive.

We study what reports of digital payment use tell us about financial inclusion. By \emph{partial inclusion} we mean that a person can use a service for one task while relying on others for a different step or lacking a way to address a failure. It is not an index or a classification of women. We draw on a survey of 423 women recruited through migrant community leaders and referrals in Colombia from December 2025 to February 2026. Recruitment aimed to reach women with international migration experience, including return migrants; the instrument did not verify every recruitment criterion. It screened for receipt of international remittances by the respondent or someone in her household over the previous year. We interpret the survey alongside written syntheses of interviews and focus groups with Venezuelan migrant women, Colombian institutional and regional expert accounts, and a comparison of five payment-system designs. The sources do not follow the same people or transactions.

Domestic digital use was common in the recruited sample. Of the 413 women with valid answers, 368 (89.1\%) reported sending or receiving money digitally within Colombia in the previous 60 days. Separately, 169 of the 423 respondents (40.0\%) reported keeping or using some incoming remittance proceeds digitally. The latter is a report from the household-screened sample, not a rate among confirmed personal recipients. The two percentages do not describe a transfer passing from one stage to the next. Among 404 scorable pairs of answers on two onboarding conditions, 100 women reported both an own-name phone line and own-name proof of address, while 141 reported neither. A phone line alone was much more common than address proof alone (Figure~\ref{fig:onboarding-conditions}). Both questions concern reported conditions; neither establishes account control or verified eligibility.

Two adjusted comparisons describe how those conditions and reported abilities are associated with the outcomes. In the remittance model, the average predicted probability is 16.35 percentage points higher at the two-condition score than at the pooled one-condition score. The model does not distinguish phone-only from proof-only at fixed covariates. In the domestic-use model, the average predicted probability is 8.15 points higher at 0.75 than at 0.50 on a score combining reported transfer ability and confidence in avoiding fraud (Figure~\ref{fig:determinants-core}). Each comparison has both exact endpoints in 11 of 12 education-by-residence strata, with an endpoint outside the observed range in one stratum. The averages retain all fitted observations. Neither comparison identifies what supplying documentation or offering training would change.

The migrant accounts give substance to these distinctions. Women describe making payments through relatives' accounts and arranging transfers abroad through intermediaries. Using someone else's account can make a purchase possible without giving the user the same authority as the holder. More direct inward transfers described in the synthesis are not a comparison group for outward journeys. Institutional accounts help explain why a document can be recognized in principle yet fail a provider's check and why opening an account differs from keeping it usable. Formal rules also identify a complaint channel for domestic Bre-B payments, not a remedy for every foreign leg \citep{banrepfaq}. The joint display (Table~\ref{tab:joint-evidence}) compares these strands without treating a shared topic as evidence about the same transaction.

Existing studies show why entry matters. \citet{ibanez2025} study eligibility for a Colombian regularization program; \citet{barboni2025} study willingness to engage with providers, and the accompanying \citet{ipaMigrantsDiagnostic} covers product ownership, wallets and institutional barriers. Research on remittance channels also distinguishes routine payment habits from how migrants send money abroad \citep{kosseVermeulen2014,martinRemittanceJourney2019}. We build on those distinctions by asking what the use indicators in this setting leave unresolved: whether women can act with or without help, whether an international recipient can use the money, and who is responsible when a transaction fails.

The framework sets out these distinctions before we describe Colombia's payment systems and the survey. We then compare the survey, migrant accounts and institutional evidence. The policy discussion identifies questions for evaluation about keeping an account usable, what a transfer costs and whether the recipient receives the full quoted amount, and what happens when a user seeks redress. The study does not estimate welfare gains or the effects of Bre-B.
